\documentclass[11pt]{article}

% The paper: MicroInfer and semantic-aware adaptive KV cache precision.
% A skeleton (#26). Sections marked TODO are to be written; the research
% notes' outline is docs_research/semantic-aware-adaptive-kv-cache.md, and
% every figure quoted must come from experiments/logs/benchmark.jsonl.

\usepackage[utf8]{inputenc}
\usepackage[T1]{fontenc}
\usepackage{amsmath}
\usepackage{booktabs}
\usepackage{hyperref}

\title{Semantic-Aware Adaptive KV Cache Precision\\
  under Video Memory Contention on Consumer GPUs}
\author{TODO}
\date{}

\begin{document}

\maketitle

\begin{abstract}
TODO.
\end{abstract}

\section{Introduction}
\label{sec:introduction}
% TODO: the problem (contention on a consumer GPU ends a session outright), the
% research questions, the contributions.

\section{Background and Related Work}
\label{sec:background}
% TODO: the KV cache and its growth; quantisation and eviction; memory
% management on consumer GPUs; related work grouped by the assumption each
% family makes that fails in a single-user local setting.

\section{Design}
\label{sec:design}
% TODO: the VRAM monitor, the attention scorer, the precision controller and
% the paged mixed-precision cache (ADR-0005, ADR-0007, ADR-0008, ADR-0011).

\section{Implementation}
\label{sec:implementation}
% TODO: an overview of MicroInfer before the boundary below.

\subsection{The From-Scratch Boundary}
\label{sec:boundary}

\textit{MicroInfer} is built from scratch in the sense this thesis depends on:
every kernel that the research contribution touches is hand-written CUDA, and no
external inference engine --- vLLM, llama.cpp or TensorRT-LLM --- is used at any
point. The phrase stops at one place, the dense matrix multiplications, which
call NVIDIA's cuBLAS library. This section states that boundary explicitly, so
that no later result is read as implying the projections were hand-written.

\paragraph{Hand-written.} The following are CUDA kernels or host code written for
this engine:
\begin{itemize}
\item the token embedding lookup;
\item RMSNorm, over fp16 activations or over the fp32 residual stream;
\item rotary position embedding (RoPE), with its angles formed in fp64 and
tabulated once per position;
\item the SwiGLU activation and the residual addition;
\item attention: causal scaled dot-product attention with an online softmax,
grouped-query attention, and the completion of keys cached without their
projection's bias --- over a contiguous cache, over pages, and over pages at a
quantised tier, dequantising each code as it is loaded;
\item the paged key-value cache: its page table, the stores that write it, and
the allocator beneath it, which uses the CUDA virtual memory management API so
that memory it frees returns to the driver;
\item quantisation and dequantisation of a page at INT8, INT4 and INT2;
\item greedy token selection over the logits.
\end{itemize}

\paragraph{cuBLAS.} The dense projections --- the query, key, value and output
projections of attention, the gate, up and down projections of the MLP, and the
language-model head --- call \texttt{cublasGemmEx}, through a single wrapper.
cuBLAS is a basic linear algebra library, not an inference engine: it multiplies
the matrices it is given and knows nothing of models, caches or precision tiers.

\paragraph{Why the boundary is here.} Hand-writing the matrix multiplications as
well was considered and rejected. Every latency figure this thesis reports would
inherit the gap between a hand-written GEMM and cuBLAS, and a reader could no
longer separate the overhead of the adaptive precision policy from the cost of a
slow matrix multiplication --- the very separation the systems claim rests on.
When the decision was taken the gap was estimated at two to five times. It was
then measured, in the study described at the end of this section: on the models'
own projections at 512 rows, a shared-memory-tiled kernel ran 8--28 times slower
than cuBLAS and a naive one 93--267 times slower, and in single-token decoding the
tiled kernel ran 5--21 times slower. The measurement strengthens the decision
rather than revising it.

\paragraph{Two precision choices.} The wrapper makes two choices on which the
per-kernel correctness figures reported for the projections depend --- a maximum
error below four and a mean below one unit in the last place of fp16, against a
float64 reference.
\begin{enumerate}
\item \textbf{fp32 accumulation, with reduced-precision reduction disallowed.}
Operands are fp16 and the computation is fp32 (\texttt{CUBLAS\_COMPUTE\_32F}).
The handle is also set to
\texttt{CUBLAS\_MATH\_DISALLOW\_REDUCED\_PRECISION\_REDUCTION}: without it,
cuBLAS may reduce the partial sums of a split-K computation in the output's
precision, which is a silent fp16 accumulation inside a call that asked for fp32.
\item \textbf{The bias folded in, so the result is rounded once.} Where a
projection carries a bias, the output rows are first seeded with it and the
multiplication adds onto them. The sum is formed in fp32 inside cuBLAS and
rounded to fp16 once, rather than rounded after the multiplication and again
after a separate addition. For the same reason the output and down projections
add straight into the fp32 residual stream, and the language-model head writes
fp32 logits: neither is rounded to fp16 at all.
\end{enumerate}

\paragraph{The GEMM study.} A naive GEMM kernel and a shared-memory-tiled one were
written and profiled against cuBLAS with Nsight Compute, as a study of where
their time goes: the naive kernel is bound by its load and store queue, the tiled
one by shared memory, and cuBLAS by the tensor cores, which neither hand-written
kernel uses. Both compute the projection correctly, to the same per-kernel bound
as the engine's. The study is a learning and benchmarking artifact and nothing
more: it is not part of the inference path, and the engine's build, its import
graph and its packaged form are each tested not to reach it.

\section{Evaluation}
\label{sec:evaluation}
% TODO: contention characterisation, the pressure detector, the task-level
% comparison of allocation strategies, and the ablations.

\section{Conclusion}
\label{sec:conclusion}
% TODO.

\end{document}
